\theory{String theory}{Can the elementary particles be understood as different vibrations of one-dimensional objects, in a framework that also includes quantum gravity?}{String theory replaces point particles by strings. Different vibrational states appear as different particles, and one closed-string state has the properties expected of a graviton. The framework uses quantum mechanics, relativity, extra dimensions, supersymmetry, branes, dualities, and conformal field theory.}{Quantum gravity, black holes, cosmology, nuclear and condensed-matter models, and pure mathematics.}{String vibrations, compactification, branes, and dualities connect particle states with quantum-gravitational geometry.}

\paragraph{The problem and history}
General relativity describes gravity and spacetime on large scales, while quantum mechanics describes microscopic phenomena. Applying ordinary quantum rules to gravity produces serious difficulties. String theory began in the late 1960s as a model of the strong force, but its unwanted features suggested a quantum theory of gravity \cite{string_ref1}.

The early bosonic theory described only bosons. Superstring theories added supersymmetry, relating bosons and fermions. Five consistent ten-dimensional superstring theories were identified: type I, type IIA, type IIB, and two heterotic theories. In the 1990s, dualities suggested that they are limits of an eleven-dimensional framework called M-theory \cite{string_ref2}.

\end{historylist}
\section*{3. Theory}
\subsection*{Core theory}
\paragraph{Strings, worldsheets, and branes}
An open string is a segment; a closed string is a loop. As a string moves, it sweeps out a two-dimensional worldsheet. Quantizing its vibrational modes produces particle states. A massless closed-string mode has spin two and is interpreted as the graviton.

A brane is a higher-dimensional object on which open strings can end. D-branes carry charges and help describe gauge fields. Compactification hides extra dimensions by making them very small. The shape of compact dimensions determines particle masses, charges, and symmetries in the lower-dimensional theory \cite{string_ref1}.

\paragraph{Dualities and correspondence}
T-duality relates a string on a circle of radius $R$ to one on a circle of radius proportional to $1/R$. Other dualities relate strong coupling to weak coupling and exchange different kinds of branes. These relations imply that apparently different theories are descriptions of one underlying structure.

The AdS/CFT correspondence, proposed in 1997, relates a gravitational string theory in an anti-de Sitter space to a conformal field theory on its boundary. It has been applied to black holes, nuclear matter, and strongly interacting condensed matter, although it is a theoretical correspondence rather than an experimentally confirmed description of nature \cite{string_ref3}.

String theory has counted microscopic states of certain black holes and has produced mathematical tools in geometry, topology, and representation theory. Compactifications on Calabi--Yau spaces and other manifolds have been used to build models resembling particle physics. However, the theory has many possible vacua, lacks a complete nonperturbative definition in all settings, and has not yet made a uniquely verified prediction about our universe \cite{string_ref4}.

\section*{4. Important theorems and results}
\begin{enumerate}[leftmargin=*,itemsep=0.35em]
\item \textbf{Graviton mode.} Closed-string quantization contains a massless spin-two state interpreted as the graviton.

\item \textbf{Supersymmetry.} Superstring theories relate bosonic and fermionic states and improve consistency.

\item \textbf{Anomaly cancellation.} Consistency restricts the allowed dimensions and gauge groups of superstring theories.

\item \textbf{Duality principle.} Different perturbative string descriptions can represent the same physics.

\item \textbf{AdS/CFT correspondence.} A gravitational theory in a bulk spacetime can be related to a conformal field theory on its boundary.

\item \textbf{Black-hole microstates.} In special supersymmetric examples, string theory counts microscopic states reproducing the Bekenstein--Hawking entropy.
\end{enumerate}
\noindent\emph{Source for these results:} \cite{string_ref1}.

\theoryapplications
\theoryconnections{string_theory}{%
\item Quantum mechanics supplies superposition and probabilities, while general relativity supplies a dynamical geometry of spacetime.
\item Replacing a point particle by a one-dimensional string changes the basic calculation from a particle path to a two-dimensional surface, whose consistency is governed by quantum field theory and conformal symmetry.
\item Extra dimensions, topology, and boundary conditions determine which vibrations are allowed and therefore which particle-like states appear \cite{string_ref1,string_ref2}.
}{%
\item String theory is useful in quantum gravity because the extended string softens some short-distance singularities and naturally includes a gravitational mode.
\item Its dualities relate apparently different physical descriptions, while branes and mirror symmetry have produced methods in geometry and topology.
\item Holographic models also connect a gravitational theory in one setting with a quantum field theory in another, giving concrete calculations for strongly interacting systems \cite{string_ref1,string_ref3,string_ref5}.
}
\section*{Glossary}
\begin{description}[leftmargin=*,style=nextline,itemsep=0.3em]
\item[\textbf{String.}] A one-dimensional extended object whose different vibrational modes correspond to different particles in string theory.
\item[\textbf{Brane.}] A higher-dimensional object on which open strings can end; D-branes carry charges and give rise to gauge fields.
\item[\textbf{Compactification.}] The process of making extra spatial dimensions very small so they are not directly observable, while their shape determines particle properties.
\item[\textbf{T-duality.}] A symmetry in string theory relating a string compactified on a circle of radius $R$ to one on a circle of radius proportional to $1/R$.
\item[\textbf{AdS/CFT correspondence.}] A conjectured duality relating a gravitational string theory in anti-de Sitter space to a conformal field theory on its boundary.
\item[\textbf{Superstring.}] A version of string theory that incorporates supersymmetry, relating bosonic and fermionic states and improving theoretical consistency.
\item[\textbf{M-theory.}] An eleven-dimensional framework that unifies the five consistent ten-dimensional superstring theories through various dualities.
\item[\textbf{Graviton.}] A massless spin-2 particle that emerges as a vibrational mode of closed strings and mediates the gravitational force.
\item[\textbf{Worldsheet.}] The two-dimensional surface swept out by a string moving through spacetime.
\item[\textbf{Supersymmetry.}] A symmetry relating bosonic and fermionic states; required for superstring consistency.
\item[\textbf{Anomaly cancellation.}] A consistency condition requiring certain quantum contributions to vanish, restricting allowed dimensions and gauge groups.
\item[\textbf{Conformal field theory.}] A quantum field theory invariant under local rescalings of the metric; the worldsheet theory of a string.
\item[\textbf{Calabi–Yau manifold.}] A compact Kähler manifold with vanishing first Chern class, used in string compactifications preserving supersymmetry.
\end{description}
\section*{7. Sample Questions}
\emph{Exercise reference:} \cite{string_ref1}. The cited edition is the source for the exercise set; consult its Exercises section for the official statement and numbering.
\begin{enumerate}[leftmargin=*,itemsep=0.9em]
\item \textbf{Exercise 1.} In bosonic string theory, the mass-squared of a closed-string state is $M^2 = \frac{2}{\alpha'}(N + \tilde{N} - 2)$, where $N$ and $\tilde{N}$ are the left- and right-moving level numbers. Find the values of $N$ and $\tilde{N}$ for the massless spin-two graviton state, and verify that it is massless. \emph{Solution.} The graviton has $N = \tilde{N} = 1$ (one oscillator excitation in each sector). Then $M^2 = \frac{2}{\alpha'}(1 + 1 - 2) = 0$. The state is massless. Its two independent polarizations give spin-two, matching the graviton.

\item \textbf{Exercise 2.} Under T-duality, a string compactified on a circle of radius $R$ is equivalent to one compactified on a circle of radius $R' = \alpha'/R$. If $\alpha' = 1$ and the original radius is $R = 3$, what is the T-dual radius? \emph{Solution.} $R' = \frac{\alpha'}{R} = \frac{1}{3}$. A string on a circle of radius $3$ is physically equivalent (in terms of the mass spectrum and interaction amplitudes) to a string on a circle of radius $\frac{1}{3}$.

\item \textbf{Exercise 3.} In the bosonic string, the ground state has $N = \tilde{N} = 0$. Compute its mass-squared and explain why it is problematic for a physical theory. \emph{Solution.} $M^2 = \frac{2}{\alpha'}(0 + 0 - 2) = -\frac{4}{\alpha'} < 0$. The ground state is a tachyon (negative mass-squared), which signals a quantum instability in the bosonic string. This is one reason the bosonic string is not a viable physical theory; superstrings resolve this by removing the tachyon through GSO projection.

\item \textbf{Exercise 4.} A closed string in bosonic string theory has left- and right-moving oscillator excitations $\alpha_{-1}^i$ and $\tilde{\alpha}_{-1}^j$ acting on the ground state, where $i, j = 1, \ldots, 24$. How many independent massless states (counting polarizations) can be formed from $N = \tilde{N} = 1$? \emph{Solution.} The state $\alpha_{-1}^i \tilde{\alpha}_{-1}^j |0\rangle$ has $24 \times 24 = 576$ components. The trace $\sum_i \alpha_{-1}^i \tilde{\alpha}_{-1}^i$ gives a scalar (the dilaton), the antisymmetric part $\alpha_{-1}^i \tilde{\alpha}_{-1}^j - \alpha_{-1}^j \tilde{\alpha}_{-1}^i$ gives $\binom{24}{2} = 276$ antisymmetric tensor states, and the remaining $576 - 1 - 276 = 299$ symmetric-traceless states are the graviton polarizations. The symmetric traceless part has $\frac{24 \times 25}{2} - 1 = 299$ components, matching.

\item \textbf{Exercise 5.} In the heterotic string, the critical dimension is $D = 26$ for the left-moving bosonic sector and $D = 10$ for the right-moving superstring sector. The difference $26 - 10 = 16$ extra dimensions are compactified. Explain how the gauge group $E_8 \times E_8$ or $\mathrm{SO}(32)$ arises from these 16 compactified dimensions. \emph{Solution.} The 16 left-moving dimensions compactified on a torus $T^{16}$ contribute 16 free left-moving bosons. Quantizing left-moving momenta on the torus gives a lattice of momenta. The requirement of modular invariance constrains this lattice to be even and self-dual in dimension 16. There are exactly two such lattices: the $E_8 \times E_8$ lattice and the $\mathrm{SO}(32)$ (or $\mathrm{D}_{16}$) lattice. These determine the gauge group of the resulting ten-dimensional theory.

\end{enumerate}
\section*{8. References}
\begin{thebibliography}{9}
\bibitem{string_ref1} J. Polchinski, \emph{String Theory, Volumes I and II}, Cambridge University Press, 1998.
\bibitem{string_ref2} B. Zwiebach, \emph{A First Course in String Theory}, Cambridge University Press, 2009.
\bibitem{string_ref3} M. B. Green, J. H. Schwarz, and E. Witten, \emph{Superstring Theory}, Cambridge University Press, 1987.
\bibitem{string_ref4} C. V. Johnson, \emph{D-Branes}, Cambridge University Press, 2003.
\bibitem{string_ref5} E. Kiritsis, \emph{String Theory in a Nutshell}, Cambridge University Press, 2007.
\end{thebibliography}
